\section{The PhD chapter: formal models and asymptotic mechanisms}

\subsection{From empirical scaling to distribution-sensitive scaling}

Empirical neural scaling laws often take the form
\[
\mathcal L(N,D,C)
\approx \mathcal L_\infty
+ a_N N^{-\alpha}
+ a_D D^{-\delta}
+ a_C C^{-\gamma},
\]
where $N$ denotes parameters, $D$ data, and $C$ compute. The lecture's central move is to treat the training distribution as an evolving object rather than a fixed background assumption. Let $P_0$ be the natural or human distribution and $P_g$ the distribution available at generation $g$. The learner is then a map
\[
\mathcal A:(P_g,T_g,\mathcal M_g)\mapsto f_g,
\]
and the generator induces a second map
\[
\mathcal G:f_g\mapsto P_{g+1}.
\]
Recursive training studies the composite dynamical system
\[
P_{g+1}=\mathcal G(\mathcal A(P_g,T_g,\mathcal M_g)).
\]
Model collapse is therefore not a single scalar event. It can mean contraction of support, altered tail exponent, increasing bias, increasing variance, worsening calibration, loss of skills, or a changed scaling exponent.

\subsection{The infinite-memory model}

Let inputs $i\in\mathbb N$ follow
\[
p_i = \frac{i^{-\beta}}{\zeta(\beta)},\qquad \beta>1,
\]
and let labels be deterministic, $y_i=f^\star(i)$. Given $T$ iid samples, the infinite-memory predictor answers correctly if $i$ was observed. Its error is
\[
E_T = \sum_{i=1}^{\infty} p_i(1-p_i)^T.
\]
For small $p_i$,
\[
(1-p_i)^T \approx e^{-Tp_i},
\]
so
\[
E_T \approx \sum_i p_i e^{-Tp_i}.
\]
A threshold analysis sets $Tp_{i^\star}\approx 1$, giving
\[
i^\star \asymp T^{1/\beta}.
\]
The remaining tail mass is
\[
\sum_{i>i^\star} i^{-\beta}
\asymp \int_{i^\star}^{\infty}x^{-\beta}\,dx
\asymp (i^\star)^{1-\beta},
\]
and therefore
\[
E_T\asymp T^{-(\beta-1)/\beta}.
\]
The exponent
\[
c=1-\frac{1}{\beta}
\]
is determined by the tail. The mechanism is occupancy, not rule learning.

\subsection{Chopped tails}

Define a synthetic distribution $Q_k$ with support restricted to $i\le k$. Evaluation remains under $P_0$. The error separates into an in-support occupancy term and missing tail mass:
\[
E_{T,k}
= \sum_{i\le k}p_i\Prob(i\text{ unseen})
+ \sum_{i>k}p_i.
\]
Under the same asymptotic approximations,
\[
E_{T,k}
\asymp T^{-(\beta-1)/\beta}
+ k^{-(\beta-1)}.
\]
Equivalently, up to constants,
\[
E_{T,k}
\asymp \min(T,k^\beta)^{-(\beta-1)/\beta}.
\]
The crossover $T\asymp k^\beta$ separates sample-limited and support-limited regimes. This structure is more informative than the generic statement ``synthetic data are worse'' because it identifies the intervention required in each regime.

\subsection{Narrowed tails}

Suppose the synthetic distribution preserves support but changes the exponent from $\beta$ to $\beta'>\beta$. This is a narrower tail. The learner is now exposed to rare events at a rate governed by $\beta'$, while evaluation may remain governed by $\beta$. Depending on the precise model, one obtains altered learning exponents and crossover behaviour rather than an immediate hard floor \citep{dohmatob2024tails}.

Temperature transformation illustrates the mechanism. If $p_i\propto i^{-\beta}$ and
\[
q_i^{(\tau)}\propto p_i^{1/\tau},
\]
then
\[
q_i^{(\tau)}\propto i^{-\beta/\tau}.
\]
For $\tau<1$, $\beta/\tau>\beta$ and the tail is narrowed. Top-$p$ sampling can instead induce an explicit cutoff.

\subsection{Multiple generations}

Let each generation use $T_0$ samples from the preceding synthetic distribution and let the final model use $T$ samples from generation $n$. In the simplified tail model, each arrow introduces another scaling term. A representative expression is
\[
E_{\mathrm{test}}^{(n)}(T)
\asymp T^{-c} + C_n T_0^{-c},
\]
with $C_n$ often growing with $n$. If $T_0\asymp T$, then
\[
E_{\mathrm{test}}^{(n)}(T)\lesssim nT^{-c}
\]
in the corresponding regime. This is not a universal exact law for LLMs; it is a structural warning that each recursive stage can contribute an additional irreducible approximation or sampling term.

For an autoregressive model,
\[
q_g(y\mid x)
=\prod_{t=1}^{L}q_g(y_t\mid x,y_{<t}).
\]
An early generated token changes the conditioning history for all later tokens. Recursive generation therefore transforms both local token probabilities and trajectory distributions.

\subsection{Mixing real and synthetic data}

Let
\[
T=T_{\Real}+T_{\AI},\qquad
\pi=\frac{T_{\Real}}{T}.
\]
In the simplified chopped-tail setting, when $T_{\Real}\ll k^\beta$,
\[
E_{\mathrm{test}}
\asymp (T_{\Real}+T_{\AI})^{-(1-1/\beta)}
+ k^{-(\beta-1)}.
\]
Synthetic data reduce occupancy error but not the missing-support term. Once sufficiently rich real data become available, the asymptotic behaviour can return to
\[
E_{\mathrm{test}}
\asymp T_{\Real}^{-(1-1/\beta)}.
\]
This motivates the interpretation of synthetic data as a bootstrap under scarcity rather than a guaranteed asymptotic replacement for real data.

\subsection{Grokking under mixed data}

Suppose synthetic examples encode an easier but biased rule and real examples encode a more general rule. Optimisation can first fit the high-density synthetic structure and only later discover the real rule. The observed curve may show:
\[
\text{rapid train fit}
\rightarrow \text{validation plateau}
\rightarrow \text{late transition}.
\]
This resembles grokking \citep{power2022grokking}. The engineering consequence is that early stopping based on a single aggregate validation metric can terminate before the transition. The research consequence is that late improvement must not be asserted as proof of a specific internal mechanism without representation-level and causal evidence.

\subsection{Kernel ridge regression}

Consider Gaussian design
\[
x\sim\mathcal N(0,\Sigma),
\qquad y=x^\top w_0+\epsilon,
\qquad \epsilon\sim\mathcal N(0,\sigma_0^2).
\]
At generation $g$, ridge regression estimates
\[
\widehat w_g
=\arg\min_w
\frac{1}{T_g}\|y_g-X_gw\|_2^2
+\lambda_g\|w\|_2^2.
\]
The closed form is
\[
\widehat w_g
=\left(X_g^\top X_g+T_g\lambda_g I\right)^{-1}X_g^\top y_g.
\]
Synthetic relabeling defines
\[
y_{g+1}=X_{g+1}\widehat w_g+\epsilon_{g+1}.
\]
Writing $e_g=\widehat w_g-w_0$ yields a recursion of the schematic form
\[
e_g=A_ge_{g-1}+b_g+\xi_g,
\]
where $A_g$ propagates inherited error, $b_g$ captures systematic regularization bias, and $\xi_g$ captures new sample and label noise. The resulting risk has contributions from clean-data estimation, inherited synthetic bias, and generation-specific variance \citep{dohmatob2024regression}.

The ridge parameter creates a bias--variance trade-off. Small $\lambda$ fits synthetic noise and teacher error; large $\lambda$ shrinks the target excessively. Hence an optimal $\lambda_g^\star$ can depend on generation depth and synthetic noise. Adaptive regularization mitigates some regimes but cannot create information absent from the labels.

\subsection{Strong model collapse}

The later strong-collapse result shows that, in a supervised regression setting, even a small constant synthetic fraction can alter asymptotic behaviour, so that increasing dataset size does not recover the clean scaling law \citep{dohmatob2025strong}. The practical lesson is not that every $1\%$ synthetic mixture is disastrous. It is that a constant synthetic fraction is not asymptotically negligible merely because it is numerically small. Its effect depends on dependence structure, model class, label bias, and scaling regime.

\subsection{Verifier-based selection}

Let a generator produce candidate pairs $(x,\widetilde y)$, with latent correctness $C\in\{0,1\}$. A verifier produces acceptance $A\in\{0,1\}$. Define
\[
\phi=\Prob(A=1\mid C=1),
\qquad
\psi=\Prob(A=1\mid C=0).
\]
The accepted-set correctness is
\[
\Prob(C=1\mid A=1)
=\frac{\phi\Prob(C=1)}{
\phi\Prob(C=1)+\psi\Prob(C=0)}.
\]
Verification is useful when it sufficiently enriches correctness and does not destroy the relevant support. The verifier can be imperfect, but its selective advantage must be measurable \citep{feng2025verification}.

There are two distinct risks:

\begin{enumerate}[leftmargin=*]
  \item \textbf{False acceptance:} bad examples enter the training set.
  \item \textbf{Selection bias:} valid but unusual examples are systematically rejected.
\end{enumerate}

A verifier can therefore prevent label collapse while creating tail collapse. Both precision and support preservation must be evaluated.

\subsection{A general decomposition for recursive synthetic learning}

A useful research abstraction is
\[
\begin{aligned}
E_g ={}& E_{\mathrm{clean}}(T_g,N_g)
+E_{\mathrm{support}}(P_0,P_g)
+E_{\mathrm{label}}(f_{g-1},f^\star)\\
&+E_{\mathrm{dependence}}(D_g)
+E_{\mathrm{verification}}(V_g)
+E_{\mathrm{shift}}(P_g,P_{\mathrm{deploy}}).
\end{aligned}
\]
Each term suggests a different intervention:

\begin{longtable}{p{0.25\textwidth}p{0.30\textwidth}p{0.35\textwidth}}
\toprule
\textbf{Error term} & \textbf{Diagnosis} & \textbf{Intervention} \\
\midrule
Clean estimation & Insufficient samples or capacity & More data, capacity, compute, better optimisation \\
Support & Missing or narrowed regimes & Real anchors, targeted simulation, diversified generation \\
Label & Biased pseudo-oracle & Independent verification, relabeling, uncertainty weighting \\
Dependence & Duplicates and correlated ancestry & Lineage-aware effective sample size \\
Verification & Weak or biased filter & Calibrate verifier, stratify acceptance, preserve tails \\
Shift & Training/deployment mismatch & Fresh data, monitoring, recalibration \\
\bottomrule
\end{longtable}

\begin{keyidea}
The phrase ``model collapse'' should be replaced, whenever possible, by a specific diagnosis: support contraction, recursive label bias, dependence inflation, verifier selection bias, or deployment shift.
\end{keyidea}
